\documentclass[UTF8,11pt,fontset=none]{ctexart}
\setCJKmainfont{Songti SC}
\setCJKsansfont{Heiti SC}
\setCJKmonofont{Songti SC}
\usepackage[a4paper,margin=2.1cm]{geometry}
\usepackage{amsmath,amssymb,booktabs,graphicx,float}
\title{Full-Space L3CR 短时高精度验证}
\author{}
\date{}
\begin{document}
\maketitle

\section{确定性目标与精度准则}
对参数向量 $\theta\in\mathbb R^n$，训练目标记为
\[
F(\theta)=\frac1{2N}\sum_{i=1}^N
\lVert f_\theta(x_i)-y_i\rVert_2^2.
\]
高精度判据直接采用
\[
r(\theta)=\lVert\nabla F(\theta)\rVert_2,
\qquad
r(\theta)\le10^{-4},\ 10^{-6},\ 10^{-8}.
\]
不作归一化。不使用验证集回退。所有路线由同一 $\theta_0$ 出发，使用 CPU、单线程和 float64。

\section{Full-space L3CR}
本实验令 active dimension 等于 $n$。第 $k$ 个三次模型为
\begin{equation}
m_k(s)=g_k^\top s+\frac12s^\top B_ks+
\frac{\sigma_k}6\lVert s\rVert_3^3,
\quad
g_k=\nabla F(\theta_k),
\quad
B_k=\nabla^2F(\theta_k).
\end{equation}
$B_k$ 由 $n$ 次 HVP 构造并对称化。PG 主迭代为
\begin{align}
y^t&=s^t-\alpha_t(g_k+B_ks^t),\\
s_i^{t+1}&=\operatorname{sign}(y_i^t)
\frac{2|y_i^t|}
{1+\sqrt{1+2\alpha_t\sigma_k|y_i^t|}}.
\end{align}
高精度精修求解同一驻点方程
\[
R_k(s)=g_k+B_ks+\frac{\sigma_k}2|s|\odot s=0,
\]
其半光滑 Newton 矩阵为
\[
J_k(s)=B_k+\sigma_k\operatorname{Diag}(|s|).
\]
只有满足
\[
\frac{\lVert R_k(s)\rVert_{3/2}}
{\max\{1,\lVert g_k\rVert_{3/2}\}}
\le10^{-12}
\]
的子问题解才允许进入外层接受检验。正则参数更新为
\[
\sigma_{k+1}=
\begin{cases}
0.5\sigma_k,&\rho_k\ge0.75,\\
\sigma_k,&0.10\le\rho_k<0.75,\\
2\sigma_k,&\rho_k<0.10\ \text{或子问题未收敛}.
\end{cases}
\]

\section{BAM 输出头的完整训练目标}
冻结主干后，每个输出通道均写成
\[
\widehat y_c=A_cw_c,
\qquad c\in\{S_1,U\},
\qquad w_c\in\mathbb R^{10}.
\]
令 $w=(w_{S_1}^\top,w_U^\top)^\top\in\mathbb R^{20}$。对全部 137 个训练案例累积
\[
G_c=A_c^\top A_c,
\qquad b_c=A_c^\top y_c,
\qquad q_c=y_c^\top y_c.
\]
于是
\begin{align}
F(w)
&=\frac1{2N}\sum_c
\left(w_c^\top G_cw_c-2b_c^\top w_c+q_c\right),\\
\nabla F(w)
&=\frac1N
\begin{bmatrix}G_{S_1}w_{S_1}-b_{S_1}\\G_Uw_U-b_U\end{bmatrix},\\
\nabla^2F(w)
&=\frac1N\operatorname{blkdiag}(G_{S_1},G_U).
\end{align}
材料体素数为 $2,159,924$。五个 Hessian 均为满秩；直接损失与 Gram 目标的最大误差为 $1.998e-15$。

\section{数值结果}
\begin{table}[H]
\centering
\caption{17 参数非线性回归：近解初始化。}
\begin{tabular}{lrrrr}
\toprule
方法 & $\lVert\nabla F\rVert_2$ & $F$ & 达到 $10^{-8}$ & AD 等价量\\
\midrule
AdamW & 1.665e-06 & 6.545e-09 & 0/10 & 201.0 \\
AdamW-0 & 1.661e-06 & 6.546e-09 & 0/10 & 201.0 \\
L-BFGS & 7.224e-07 & 5.184e-10 & 0/10 & 490.1 \\
Full-space L3CR & 2.698e-06 & 7.448e-11 & 1/10 & 474.1 \\
\bottomrule
\end{tabular}
\end{table}

\begin{table}[H]
\centering
\caption{17 参数非线性回归：远解负对照。}
\begin{tabular}{lrrrr}
\toprule
方法 & $\lVert\nabla F\rVert_2$ & $F$ & 达到 $10^{-8}$ & AD 等价量\\
\midrule
AdamW & 1.742e-04 & 1.439e-05 & 0/10 & 201.0 \\
AdamW-0 & 1.742e-04 & 1.439e-05 & 0/10 & 201.0 \\
L-BFGS & 6.959e-07 & 4.790e-08 & 0/10 & 451.7 \\
Full-space L3CR & 3.067e-04 & 1.480e-05 & 0/10 & 495.0 \\
\bottomrule
\end{tabular}
\end{table}

\begin{table}[H]
\centering
\caption{BAM-KAN 20 参数输出头。}
\begin{tabular}{lrrrr}
\toprule
方法 & $\lVert\nabla F\rVert_2$ & $F$ & 达到 $10^{-8}$ & AD 等价量\\
\midrule
AdamW & 3.836e-05 & 1.497e-01 & 0/5 & 201.0 \\
AdamW-0 & 3.836e-05 & 1.497e-01 & 0/5 & 201.0 \\
L-BFGS & 4.741e-08 & 1.497e-01 & 0/5 & 502.8 \\
Full-space L3CR & 2.318e-09 & 1.497e-01 & 5/5 & 71.4 \\
\bottomrule
\end{tabular}
\end{table}

\begin{figure}[H]
\centering
\includegraphics[width=0.94\linewidth]{high_accuracy_convergence.pdf}
\caption{原始梯度范数与反向传播等价量。虚线为 $10^{-8}$。}
\end{figure}

\section{结论边界}
总实验时间为 $116.38\,\mathrm{s}$。BAM 输出头上，L3CR 在 $5/5$ 个 checkpoint 达到 $10^{-8}$；其余三条路线均为 $0/5$。该结果支持20维满秩二次输出头校正中的条件性高精度优势。

17参数非线性近解实验中，L3CR仅在 $1/10$ 个 seed 达到 $10^{-8}$；远解负对照为 $0/10$。故不能推出一般非线性神经网络优势，也不能外推至完整 86,653 参数 BAM-KAN。
\end{document}
